\documentclass[a4paper]{article}
% Español (México) con XeLaTeX: fontspec para fuentes Unicode, polyglossia
% para la división silábica y los nombres del documento en la variante mexicana.
\usepackage{fontspec}
\usepackage{polyglossia}
\setdefaultlanguage[variant=mexican]{spanish}
% Los nombres chinos van como «pinyin (original)»: xeCJK da a los caracteres
% Han su propia fuente; sin ella XeLaTeX los omite con un aviso.
% FandolSong no cubre algunos caracteres de nombres (U+73FA); WenQuanYi Zen Hei sí.
\usepackage[AutoFallBack=true]{xeCJK}
\setCJKmainfont{FandolSong-Regular.otf}
\setCJKfallbackfamilyfont{\CJKrmdefault}{WenQuanYi Zen Hei}
% polyglossia no traduce los nombres de listings.
\AtBeginDocument{\providecommand{\lstlistingname}{}\renewcommand{\lstlistingname}{Listado}%
  \providecommand{\lstlistlistingname}{}\renewcommand{\lstlistlistingname}{Índice de listados}}
\usepackage{amsmath, amssymb}
\usepackage{graphicx}
\usepackage[margin=2.5cm]{geometry}
\usepackage[most]{tcolorbox}
\usepackage{etoolbox}
\usepackage{listings}
\usepackage{hyperref}
\usepackage{booktabs}
\usepackage{float}
\newtcolorbox{knowledgebox}[1]{enhanced, colback=blue!5!white, colframe=blue!75!black, colbacktitle=blue!75!black, coltitle=white, fonttitle=\bfseries, title={#1}, attach boxed title to top left={yshift=-2mm, xshift=2mm}, boxrule=1pt, sharp corners}
\newtcolorbox{importantbox}[1]{enhanced, colback=yellow!10!white, colframe=yellow!80!black, colbacktitle=yellow!80!black, coltitle=black, fonttitle=\bfseries, title={#1}, sharp corners}
\newtcolorbox{warningbox}[1]{enhanced, colback=red!5!white, colframe=red!75!black, colbacktitle=red!75!black, coltitle=white, fonttitle=\bfseries, title={#1}, sharp corners}
\lstset{language=Python, basicstyle=\ttfamily\small, keywordstyle=\color{blue}, stringstyle=\color{red!60!black}, commentstyle=\color{green!60!black}, breaklines=true, frame=single, numbers=left, numberstyle=\tiny\color{gray}, captionpos=b, extendedchars=true}
\newcommand{\notetitle}{Instalación y desarrollo del entorno veRL basado en Docker}
\newcommand{\noteauthors}{Elaboradas a partir de materiales públicos del curso}
\newcommand{\notedate}{2025}
\newcommand{\videochannel}{Wudaokou Nashi (五道口纳什)}
\newcommand{\videopublishdate}{2025}
\newcommand{\videoduration}{aproximadamente 20 minutos}
\newcommand{\videourl}{}
\newcommand{\videocoverpath}{cover.jpg}
\newcommand{\repourl}{https://github.com/hqhq1025/ai-course-notes}

% Footer with repo info
\usepackage{fancyhdr}
\pagestyle{fancy}
\fancyhf{}
\fancyfoot[C]{\thepage}
\fancyfoot[R]{\footnotesize\color{gray}\href{\repourl}{AI Course Notes}}
\renewcommand{\headrulewidth}{0pt}
\renewcommand{\footrulewidth}{0pt}

\begin{document}
\begin{titlepage}\centering
{\Large Notas del curso\par}\vspace{1.2cm}{\huge\bfseries \notetitle\par}\vspace{0.8cm}{\large \noteauthors\par}\vspace{0.3cm}{\large \notedate\par}\vspace{1.2cm}
\ifdefempty{\videocoverpath}{}{\includegraphics[width=0.82\textwidth,height=0.45\textheight,keepaspectratio]{\videocoverpath}\par}
\vfill
\begin{tcolorbox}[width=0.9\textwidth, colback=black!2!white, colframe=black!60, sharp corners]
\textbf{Autor o canal del video}: \videochannel\par\textbf{Fecha de publicación}: \videopublishdate\par\textbf{Duración del video}: \videoduration\par
\ifdefempty{\videourl}{}{\textbf{Enlace del video}: \href{\videourl}{\nolinkurl{\videourl}}\par}\par
\textbf{Página del proyecto}: \href{\repourl}{\nolinkurl{\repourl}}\par
\end{tcolorbox}
\end{titlepage}
\tableofcontents\newpage

\section{Introducción}
Esta entrega presenta la instalación de veRL basada en Docker, para lograr una configuración del entorno sin complicaciones.

\begin{knowledgebox}{¿Por qué usar Docker?}
El método tradicional de pip install a menudo presenta conflictos de dependencias de versión entre Flash Attention y PyTorch. Docker proporciona un entorno aislado y reproducible que evita estos problemas.
\end{knowledgebox}

\section{Conceptos básicos de Docker}
\subsection{Imágenes y contenedores}
\begin{itemize}
  \item \textbf{Image} (imagen): plantilla de solo lectura que contiene el sistema operativo y el entorno de software
  \item \textbf{Container} (contenedor): instancia en ejecución de una imagen, de lectura y escritura
  \item La relación es similar a la de «clase» y «objeto»
\end{itemize}

\subsection{La máquina anfitriona y el contenedor}
Entorno de desarrollo: Mac (terminal local) $\to$ servidor GPU remoto $\to$ veRL se ejecuta dentro de un contenedor Docker. Se conecta al contenedor remoto mediante VSCode/Cursor para lograr una experiencia de desarrollo local.

\section{Comando principal}
\texttt{docker run -it --rm -v <host>:<container> -w <workdir> <image>:<tag> <cmd>}
\begin{itemize}
  \item \texttt{-it}: terminal interactiva
  \item \texttt{--rm}: elimina el contenedor al salir
  \item \texttt{-v}: montaje de disco (mapea un directorio del host al contenedor)
  \item \texttt{-w}: establece el directorio de trabajo
\end{itemize}

\subsection{Resumen de la sección}
Docker es la mejor práctica para administrar entornos complejos de ML. veRL ofrece de manera oficial una imagen de Docker con la que se puede configurar rápidamente un entorno de entrenamiento.

\newpage
\section{Recomendación práctica: primero estabilizar el entorno, después hablar de eficiencia de entrenamiento}
Para quien empieza, el mayor valor de Docker no es «ser más avanzado», sino hacer que el entorno de entrenamiento sea reproducible. Sobre todo en un flujo de trabajo de GPU remota con IDE local, basta con definir con claridad la imagen, los montajes y el directorio de trabajo para eliminar de antemano la mayoría de los problemas de entorno.

\begin{knowledgebox}{Flujo de trabajo mínimo recomendado}
\begin{itemize}
  \item Primero usar la imagen oficial para verificar que las dependencias estén completas
  \item Después usar un directorio montado para guardar el código y las salidas de los experimentos
  \item Por último conectar VSCode Remote para el desarrollo cotidiano
\end{itemize}
\end{knowledgebox}

\subsection{Resumen de la sección}
Docker no es una herramienta accesoria, sino infraestructura de la ingeniería de aprendizaje automático. Resolver primero el problema de la reproducibilidad del entorno y después optimizar la velocidad de entrenamiento suele ser más conveniente.

\newpage
\section{Síntesis y ampliación}
\begin{enumerate}
  \item Docker resolvió el problema de conflictos de dependencias del entorno
  \item El flujo de trabajo de imagen $\to$ contenedor $\to$ montaje de disco
  \item VSCode Remote se conecta al contenedor para lograr una experiencia de desarrollo local
\end{enumerate}
\end{document}

